\section{Análisis de Resultados}

\subsection{Comparación entre la respuesta teórica y la real}

Al observar los resultados gráficos se puede notar que para el observador identidad la posición oscila tres veces antes de estabilizarse, ambas curvas apuntan al mismo comportamiento. Los cambios de amplitudes son similares y los picos cambian en los mismos instantes de tiempo en ambas gráficas. Vemos que para ambos casos el error en estado estacionario es prácticamente 0. En los momentos en que el error estacionario no es cero en el prototipo real cuando debería serlo se debe a perturbaciones externas como el viento.

Al evaluar las respuestas del prototipo físico con los distintos observadores se nota que la respuesta de posición para el observador identidad es subamortiguada mientras que la del de orden reducido es más parecida a una respuesta críticamente amortiguada o sobreamortiguada ya que no tiene sobrepico. Por otro lado las respuestas del observador de orden completo y del funcional lineal parecen tener un pequeño sobreimpulso, sin embargo esto puede ser debido a la perturbación del viento. Por otro lado para el ángulo del péndulo vemos que se comportan prácticamente de la misma manera.

\subsection{Comparación entre los controladores}

Se observa claramente en las gráficas que la respuesta de posición del PID tiene un comportamiento subamortiguado mientras que para los tres controladores por realimentación de estados es prácticamente sobreamortiguada. Vemos que en estado estable el PID mantiene pequeñas oscilaciones mientras que los demás se mantienen prácticamente con error 0. Esto se debe principalmente a que el PID no puede adaptarse de forma tan óptima a los modos acoplados del lazo cerrado como sí lo hacen los de variables de estado. Al observar el ángulo se nota que el PID llega a ángulos mayores llegando a registrar picos de hasta 5 grados mientras que los demás esquemas se mantuvieron alrededor de 1 grado.

En la velocidad del carro se observa que el PID llega a valores más altos de hasta 8 mientras que los demás tienen su pico mayor en 3. Se notó además que en estado estacionario la velocidad del PID se encuentra muy ruidosa mientras que los otros tienen una velocidad que oscila muy cerca de cero en estado estacionario. Esta puede ser la razón por la cual los controladores de variables de estado son más robustos que el PID, estos realizan más trabajo interno continuo para mantener el ángulo en 0. En la aceleración ocurre algo parecido que con la velocidad. Para el PID la aceleración es muy ruidosa en estado estable mientras que en los demás se mantiene en un rango estrecho de entre -5 y 5. Se observó también que la aceleración en el funcional lineal llega a 5 durante la prueba de cambio de referencia mientras que en estado estacionario se reduce a valores cercanos a cero.

Por otro lado el de orden reducido maneja valores de aceleración menores que el resto en estado estacionario tanto en el cambio de referencia como en estado estacionario. Esto podría indicar que es el más eficiente energéticamente de los tres. Cabe mencionar que los límites de saturación de velocidad y aceleración son los mismos para todos los controladores.

Las figuras muestran los resultados de las pruebas de condiciones iniciales para ángulos de 10°, 15°, 25° y 30° respectivamente. En la prueba de 10° se obtuvo que el PID tuvo la mayor sobreelongación luego le sigue el funcional lineal mientras que el identidad y el de orden reducido tuvieron una sobreelongación menor. Tanto el PID, identidad y el funcional lineal oscilaron al menos 2 veces antes de estabilizarse, mientras que el control con observador de orden reducido se estabilizó a través de la primera sobreelongación.

Para la prueba de 15° el PID no logró recuperarse y se cayó mientras que los otros lograron estabilizarse. En esta ocasión el que mostró una respuesta más suave fue el observador identidad que osciló una sola vez mientras que los otros oscilaron al menos dos veces. En las pruebas de 25° y 30° no se probó el PID ya que se probó antes que no es capaz de recuperarse de ángulos mayores a 15°. En esta ocasión los tres controladores por variables de estado se comportaron de manera similar lo que parece indicar que el comportamiento sobreamortiguado depende de las condiciones del sistema y de las perturbaciones al momento de hacer las pruebas.

En la prueba con el ventilador se notó que el PID no fue capaz de soportar la perturbación continua y terminó cayendo, mientras que los de variables de estado se mantuvieron estables. Entre estos últimos el funcional lineal es el que más esfuerzo hizo para mantener la estabilidad seguido por el de orden reducido y por último el identidad.

Durante la calibración del PID también se observó que aumentar la ganancia proporcional hacía que el carro se desplazara bruscamente generando vibraciones mecánicas. Esto limitó en gran medida el ajuste del PID. Esta dificultad parece ser debida al ruido en la lectura del ángulo del encoder.
